% Generated draft: requires mathematical and source review.
\documentclass[11pt]{article}
\usepackage[margin=25mm]{geometry}
\usepackage{fontspec}
\setmainfont{Latin Modern Roman}
\usepackage{amsmath,amssymb,mathtools,unicode-math}
\setmathfont{Latin Modern Math}
\usepackage{enumitem,needspace,xurl,hyperref}
\setlength{\parindent}{0pt}
\setlength{\parskip}{5pt}
\setlength{\emergencystretch}{3em}
\newcommand{\N}{\mathbb{N}}
\newcommand{\Z}{\mathbb{Z}}
\newcommand{\Q}{\mathbb{Q}}
\newcommand{\R}{\mathbb{R}}
\newcommand{\C}{\mathbb{C}}
\newcommand{\F}{\mathbb{F}}
\newcommand{\A}{\mathbb{A}}
\newcommand{\PP}{\mathbb{P}}
\newcommand{\E}{\mathbb{E}}
\newcommand{\eps}{\varepsilon}
\newcommand{\dd}{\,\mathrm{d}}
\begin{document}
\section*{UnsolvedMath ID: 30000897}
\textbf{Source problem code:} catalogue entry
\subsection{Definitions and mathematical statement}
Work in the Euclidean plane \(\R^2\). A direction \(u\) is a one-dimensional linear subspace, so opposite vector representatives specify the same direction. Let \(\mathcal L_u\) be the set of all affine lines parallel to \(u\). For a finite set \(A\subset\R^2\), its discrete X-ray in direction \(u\) is the function
\[
X_uA:\mathcal L_u\to\Z_{\geq0},\qquad
(X_uA)(\ell)=|A\cap\ell|.
\]
The locations of the lines are part of the data. Given \(k\geq1\) distinct directions \(u_1,\ldots,u_k\), call \(A\) reconstructible if, for every finite set \(B\subset\R^2\),
\[
\bigl(X_{u_i}B=X_{u_i}A\text{ for all }1\leq i\leq k\bigr)
\quad\Longrightarrow\quad B=A.
\]
Thus reconstruction is exact in the fixed plane, not only up to translation or other congruence. There is no lattice, convexity, generic-position, or prescribed-support restriction in the supplied statement.

The source defines
\[
f(u_1,\ldots,u_k)=\max\{n:\text{every }A\subset\R^2\text{ with }|A|=n
\text{ is reconstructible from these directions}\},
\]
and
\[
F(k)=\max_{u_1,\ldots,u_k\text{ distinct}} f(u_1,\ldots,u_k).
\]
Determine \(F(k)\), or improve estimates for it as a function of \(k\). The same chosen directions must work for every set of the stated size; they are chosen before the unknown set, rather than adapted to that set. The record does not state whether zero is allowed in the first maximum when no positive size works, or give separate endpoint conventions; these require confirmation for small \(k\). The maxima are retained as written rather than replaced by an unproved claim about their existence or a different optimization problem.
\subsection{Short English statement}
Choose \(k\) directions to maximize the size for which every finite planar point set is uniquely determined by its line-count X-rays. Determine that optimal guaranteed size \(F(k)\), or obtain better bounds.
\subsection{Sources}
\begin{itemize}
\item Local UnsolvedMath registry, permanent ID \(30000897\), source code catalogue entry, ``Point-Set Reconstruction from Directional X-Rays.'' No individual proposer is supplied.
\item \emph{Combinatorics} (2008), Oberwolfach Reports, DOI 10.4171/owr/2008/01, pp. 5--78: \url{https://doi.org/10.4171/owr/2008/01}.
\item Supplied reference, ``discrete X-ray reconstruction paper'': \url{https://doi.org/10.1016/j.endm.2007.07.069}.
\item Supplied literature triage, checked 2026-08-21: status ``open.'' It reports exponential upper and lower bounds in the defining paper, without an exact general answer, and calls the reconstruction problem unresolved. No constants or ranges for those bounds are supplied. This dated assessment derives from supplied OWR curation; time-sensitive status and cited claims require independent verification.
\end{itemize}
\end{document}
